% Intended LaTeX compiler: pdflatex
\documentclass[a4paper,12pt]{article}
\usepackage[utf8]{inputenc}
\usepackage[T1]{fontenc}
\usepackage{graphicx}
\usepackage{longtable}
\usepackage{wrapfig}
\usepackage{rotating}
\usepackage[normalem]{ulem}
\usepackage{amsmath}
\usepackage{amssymb}
\usepackage{capt-of}
\usepackage{hyperref}
\usepackage{\string~/.emacs.d/common}
\author{mahmood}
\date{\today}
\title{algorithms 2 homework 2}
\hypersetup{
 pdfauthor={mahmood},
 pdftitle={algorithms 2 homework 2},
 pdfkeywords={},
 pdfsubject={},
 pdfcreator={Emacs 28.2 (Org mode 9.5.5)}, 
 pdflang={English}}
\begin{document}

\maketitle
\begin{question}
improve the algorithm \(\textsc{Dijkstra}\) so that it knows for every node \(v\) how many edges are in one of the shortest paths from the root node to \(v\)\\
show the correctness of the algorithm and its \href{20221130014441-time_complexity.org}{time complexity}\\
\begin{answer}
\includesvg{/home/mahmooz/brain/out/0vBadYH}\\

the time complexity is \(O(|E|\log|V|)\)\\
\end{answer}
\end{question}
\begin{question}
improve the algorithm \(\textsc{BFS}\) so that it knows for every node \(v\) how many short paths there is from the root node to \(v\)\\
show the correctness of the algorithm and its time complexity\\
\begin{answer}
\includesvg{/home/mahmooz/brain/out/UpgeA7b}\\

the time complexity of the algorithm is still \(\Theta(|E|+|V|)\)\\
\end{answer}
\end{question}
\begin{question}
given a directed and non-negatively weighted graph \(G=(V,E,W)\) that represents a system of roads and junctions and the distances between them. given a set of junctions \(F \subseteq V\) that contain gas stations. a car can carry fuel that allows it to travel \(K\) kilometers. given a starting junction \(u\) and a destination junction \(v\), propose an algorithm that answers the question: can the car travel from the source junction to the destination junction (assume that the car starts with a full fuel container).\\
\begin{answer}
we use a modified version of \href{20230412172450-dijkstra_s_algorithm.org}{dijkstra's algorithm}\\
\includesvg{/home/mahmooz/brain/out/7y8kdId}\\
\end{answer}
\end{question}
\begin{question}
\begin{subquestion}
show an example of a non-\href{20230412150436-weighted_graph.org}{weighted graph} in which a \href{20230415232312-shortest_path_tree.org}{distance tree} relative to a specific vertex isnt the same for another vertex\\
\begin{answer}
pick a triangle, connect a vertex to one of its edges, this would result in the other 2 vertices having a distance of 2 from that new vertex yet the vertex we connected the new vertex to only has a distance of 1 from all the other vertices including the new one so the distance trees must differ\\
\end{answer}
\end{subquestion}
\begin{subquestion}
prove that if a specific \href{20230415231645-spanning_tree.org}{spanning tree} of a non-weighted graph is the distance tree relative to every vertex in the graph then the graph is a \href{20230323170647-tree.org}{tree}\\
\begin{answer}
i couldnt solve this\\
\end{answer}
\end{subquestion}
\end{question}
\begin{question}
given a directed graph \(G=(V,E)\), and given \(s,t \in V, e_1 \in E\), propose an efficient algorithm that finds the shortest path that goes from \(s\) to \(t\) and goes through the edge \(e_1\)\\
explain why the algorithm works correctly, analyze the complexity of the algorithm\\
\begin{answer}
couldnt solve this either\\
\end{answer}
\end{question}
\end{document}